% فصل پنجم: نتیجه‌گیری و پیشنهادها
\chapter{نتیجه‌گیری و پیشنهادها}
\label{ch:conclusion}

در این فصل، ابتدا خلاصه‌ای از آنچه در این پایان‌نامه انجام شد ارائه می‌شود،
سپس یافته‌های اصلی و نتیجه‌گیری نهایی بیان می‌گردد و در پایان، پیشنهادهایی
برای ادامه پژوهش در این زمینه مطرح می‌شود.

\section{خلاصه پژوهش}
\label{sec:summary}

این پایان‌نامه به بررسی چارچوب شبکه‌های عصبی آگاه از فیزیک برای حل مسائل
مستقیم و معکوس معادلات دیفرانسیل جزئی غیرخطی اختصاص داشت. در فصل اول،
موضوع، مسئله، اهداف و سؤال‌های پژوهش تعریف شد. در فصل دوم، مبانی نظری لازم
شامل معادلات دیفرانسیل جزئی، روش‌های عددی کلاسیک، شبکه‌های عصبی و آموزش
آن‌ها، مشتق‌گیری خودکار و پیشینه پژوهش در زمینه یادگیری ماشین آگاه از
فیزیک مرور گردید. در فصل سوم، چارچوب اصلی پژوهش یعنی تعریف شبکه باقیمانده،
طراحی تابع هزینه ترکیبی، مدل‌های زمان‌پیوسته و زمان‌گسسته و فرمول‌بندی
مسئله معکوس تشریح شد. در فصل چهارم، مدل زمان‌پیوسته برای معادله برگرز
یک‌بعدی با زبان پایتون و کتابخانه جکس پیاده‌سازی شد؛ یک حل‌گر مرجع مستقل
بر پایه روش طیفی فوریه و گام‌زنی \lr{ETDRK4} توسعه یافت و با دو آزمون
مستقل راستی‌آزمایی گردید؛ و نتایج حل مستقیم و کشف پارامترها گزارش و تحلیل
شد.

\section{یافته‌ها و نتیجه‌گیری}
\label{sec:findings}

یافته‌های اصلی این پژوهش را می‌توان به‌شرح زیر خلاصه کرد:

\begin{enumerate}
\item مدل زمان‌پیوسته شبکه آگاه از فیزیک توانست جواب معادله برگرز را در کل
دامنه زمانی ـ مکانی، تنها با ۲۰۰ نقطه داده اولیه و مرزی بازسازی کند و به
خطای نسبی \FwdRelErr\ در نرم $L_2$ نسبت به حل مرجع دست یابد. این نتیجه،
کارایی داده‌ای بالای این چارچوب را تأیید می‌کند.

\item خطای جواب تقریباً به‌طور کامل در لایه باریک جبهه موج متمرکز بود و در
بقیه دامنه ناچیز به دست آمد. تمرکز نیمی از نقاط هم‌مکانی در ناحیه جبهه،
نقش مهمی در دستیابی به این نتیجه داشت و نشان داد که توزیع نقاط هم‌مکانی
باید با فیزیک مسئله (محل گرادیان‌های تند) هماهنگ باشد.

\item در مسئله معکوس، ضریب انتقال ($\lambda_1$) با خطای \LamIErrClean\ و
ضریب پخش ($\lambda_2$) با خطای \LamIIErrClean\ شناسایی شد. خطای بیشتر ضریب
پخش، ریشه در کوچکی مقدار آن و موضعی بودن اثرش در لایه جبهه دارد و با
یافته‌های مقاله مرجع سازگار است.

\item افزودن یک درصد نویز به داده‌های مسئله معکوس، دقت شناسایی را به‌طور
چشمگیری تخریب نکرد (خطای \LamIErrNoisy\ برای $\lambda_1$ و \LamIIErrNoisy\
برای $\lambda_2$)، که پایداری قابل‌قبول روش در برابر نویز را نشان می‌دهد.

\item مقایسه با مقاله مرجع نشان داد که با امکانات محدود (پردازنده مرکزی و
بودجه محاسباتی کمتر) می‌توان رفتار کیفی روش را به‌طور کامل بازتولید کرد،
هرچند رسیدن به دقت عددی گزارش‌شده در مقاله نیازمند توان محاسباتی بیشتر و
تنظیم دقیق‌تر ابرپارامترهاست.
\end{enumerate}

در مجموع می‌توان نتیجه گرفت که شبکه‌های عصبی آگاه از فیزیک، چارچوبی ساده،
کلی و کم‌داده برای ترکیب دانش فیزیکی با یادگیری ماشین ارائه می‌دهند و
به‌ویژه در مسائلی که داده تجربی کمیاب است یا پارامترهای مدل مجهول‌اند،
می‌توانند مکمل ارزشمندی برای روش‌های عددی کلاسیک باشند. در عین حال، این
روش‌ها هنوز با چالش‌هایی مانند تفکیک لایه‌های تند جواب، هزینه بالای آموزش
نسبت به یک اجرای حل‌گر کلاسیک و نبود تضمین‌های نظری همگرایی مواجه‌اند و
نباید آن‌ها را جایگزینی بی‌قیدوشرط برای روش‌های کلاسیک دانست.

\section{پیشنهادها برای پژوهش‌های آینده}
\label{sec:future}

در ادامه این پژوهش، موضوع‌های زیر پیشنهاد می‌شود:

\begin{enumerate}
\item پیاده‌سازی مدل زمان‌گسسته با طرح‌های رانگ ـ کوتای ضمنی و مقایسه
کارایی و دقت آن با مدل زمان‌پیوسته روی همان مسئله برگرز؛

\item اعمال روش به معادلات دیگر مانند معادله گرما، معادله شرودینگر غیرخطی
یا معادلات ناویر ـ استوکس در هندسه‌های ساده؛

\item بررسی راهبردهای بهبود دقت در ناحیه جبهه، از جمله وزن‌دهی تطبیقی
جملات تابع هزینه و نمونه‌برداری تطبیقی نقاط هم‌مکانی در حین آموزش؛

\item مطالعه نظام‌مند اثر معماری شبکه (عمق و پهنا)، تعداد نقاط داده و
هم‌مکانی و راهبرد بهینه‌سازی بر دقت نهایی؛

\item ترکیب چارچوب آگاه از فیزیک با روش‌های برآورد عدم قطعیت، مانند
شبکه‌های بیزی یا روش‌های مبتنی بر اجماع، برای کاربردهای حساس به اطمینان؛

\item بررسی مسائل دوبعدی و ارزیابی مقیاس‌پذیری روش با افزایش ابعاد مسئله.
\end{enumerate}
